\newpage
\section{Foundation Model}

This chapter describes background architecture. The experiments use the TimesFM 2.5 checkpoint under frozen-weight inference. Architectural details from the original TimesFM paper are historical background and should not be assumed unchanged in version 2.5; the executed configuration is specified in Section~\ref{sec:protocol}.\cite{timesfm_paper,timesfm_card}

% Is this even necessary? Idk, i will keep it brief for now to not spend too much on it if it isnt
\subsection {Transformer Architecture}

The Transformer is the most widely used architecture in present-day machine learning. It was initially invented for NLP tasks; however, since its inception, it has been adopted by most branches of the field, proving superior in most tasks. Its main advantage is the ability to process large sequences of data in parallel. This improves training speed and allows for a high degree of scalability, which in turn highly increases the possible performance. It is based on the self-attention mechanism that makes the model learn contextual relationships between all the elements in an input, which is well fitting for a problem such as time series forecasting. The original Transformer is an encoder-decoder model, however many of the new ideas that build on it switch to an encoder-only or decoder-only approach. It uses tokenized input sequences for training for a number of key reasons, including the use of embeddings, the creation of a finite token vocabulary, and ease of incorporating positional information into tokenized data.

\subsection{TimesFM}

TimesFM is a decoder-only approach to the problem of time series forecasting. The goal of the model is to generate a sequence of future values based on the given context. The model is designed to handle variable context lengths and prediction horizons, without being tied to a specific use case. As a foundation model, TimesFM aims to outperform or achieve comparable zero-shot results to other state-of-the-art models’ performance on previously seen data. In principle, the model uses a GPT-like decoder-only architecture, training on time series data analogous to tokens used in large language models. However, it introduces novel improvements to allow for better generalized learning on data of variable domain, granularity, context, and horizon lengths, designed with consideration to the specificity of time series data.

\subsubsection{Input patching and embedding}

TimesFM does not operate on individual time steps. Instead, similar to other recently developed long-horizon forecasting models,  it splits an entire sequence of data into fixed-length contiguous non-overlapping patches. This process is similar to the tokenization of text data and it makes the input sequence length shorter by a factor of the patch length, which greatly increases inference speed and improves forecasting performance. 
For a context window of $ N $ time steps, $ \frac{N}{\mathit{input\_patch\_len}}$ patches are created, with padding and a mask supplied if $ N $ is not an exact multiple. Every patch is then passed by an input residual block, a multilayer perceptron with one hidden layer and a skip connection. This transforms the patch into an embedding of dimension $ \mathit{d_{model}} $, which represents the model's latent dimensionality. In addition to the input data, the residual block is also supplied with a binary padding mask. For any data point $ y $, a corresponding mask element $ m $ signifies whether the input should be read and interpreted by the model or not. This allows the model to omit training on any gaps in data or purposely hidden parts of the input. The padding is primarily used for random patch masking necessary for the model to learn on different context lengths, which will be further explained below. After the residual block, a positional encoding is added to the processed embedding. This ensures the model properly recognizes the temporal position of each element in the input sequence. TimesFM uses sinusoidal positional encodings akin to the original Transformer. The final input is provided to decoder layers described in a further subsection.
The entire input processing can be denoted as follows:
\\\textbf{Throughout the whole process.}
\begin{itemize}
    \item $ L $ --- length of the input time series,
    \item $ p $ --- patch length,
    \item $ N = \lceil\frac{L}{p}\rceil$ --- number of patches,
    \item $ y_{1:L} = (y_1,\ y_2,\ \ldots,\ y_L)$ --- raw input values,
    \item $ m_{1:L} = \left( m_1,\ m_2,\ \ldots,\ m_L \right) \in \{0,\ 1\}^L $ --- binary mask.
\end{itemize}
\begin{enumerate}
    \item \textbf{Partition into patches.}\\
        For $ i = 1,\ 2,\ \ldots, \ N $, the $ i $--th patch of length $ p $ is defined by
        \[ \tilde{y}_i = \left( y_{p(i-1)+1}:y_{pi} \in \mathbb{R}^p \right),\]
        similarly, the mask vector over that patch is defined by
        \[ \tilde{m}_i=\left( m_{p(i-1)+1}:m_{pi} \in \left\{ 0,\ 1 \right\}^p \right).\]
    \item \textbf{Patch masking.}\\
        Masking each element in the input data based on the binary padding can be denoted as
        \[ x_i = \tilde{y}_i\ \odot\ (1-\tilde{m}_i),\]
        where $\odot$ is the element-wise product.
    \item \textbf{Residual block embedding.}\\
        Each masked patch is processed into an embedding with an MLP defined as
        \[MLP(x) = W_2 \sigma(W_1x+b_1) + b_2,\]
        where
        \begin{itemize}
            \item $ h $ -- hidden size,
            \item $ W_1 \in \mathbb{R}^{h\times p},\ b_1\in \mathbb{R}^h $,
            \item $ W_2 \in \mathbb{R}^{d_{model}\times h},\ b_2 \in \mathbb{R}^{d_{model}}, $
            \item $ \sigma $ --- point-wise nonlinearity, e.g. ReLU. % should i expand on ReLU?
        \end{itemize}
        With a learned skip projection $W_{\mathrm{skip}}\in\mathbb{R}^{d_{model}\times p}$, the residual block then outputs
        \[e_i = W_{\mathrm{skip}}x_i + MLP(x_i) \in \mathbb{R}^{d_{model}}.\]
    \item \textbf{Positional encoding addition.}\\
        For $ j \in \{0,\ 1, \ldots\} $, positional encoding is defined by
        \[PE(j)_{2k}=sin\left(j / 10000^{\frac{2k}{d_{model}}}\right),\]
        \[PE(j)_{2k+1}=cos\left(j / 10000^{\frac{2k}{d_{model}}}\right).\]
        The final token provided to the decoder can be denoted as
        \[t_i = e_i+PE(i).\]
\end{enumerate}

\subsubsection{Handling variable context and horizon lengths}

A fundamental requirement for a foundation time series zero-shot forecasting model is the ability to handle historical contexts and predict the future behavior of sequences of lengths indeterminate at training time. TimesFM addresses this in its architectural design and pretraining process so that during inference it can accept any context length up to a set limit and produce forecasts of arbitrary horizon. It is designed to work autoregressively; thus, it iteratively feeds its own outputs back as input data to create further predictions up to a requested forecast length. It is also trained on data of varied context length, granularity, and origin. Furthermore, it uses special randomized masking during training to prepare the model for receiving data of any length, from 1 to maximum context length. Both of these procedures are described further below.

\subsubsection{Randomized context masking}

A model trained on data of a predetermined sequence length is bound to specialize in forecasting for that specific length only. Even despite dividing the data into multiple patches, such a model might only be able to make correct predictions for horizons that are multiples of the patch length. In order to generalize effectively across every possible context length up to a given maximum, TimesFM applies a randomized masking strategy to ensure the model sees examples of data of all possible sizes during training. During the input preprocessing, a mask of a randomly chosen length is applied to the beginning of every input. The mask spans across part of the first patch of data. This means the model cannot use that data to train on, effectively presenting it with data of a shorter context than the full sequence. Over the course of training, the mask length is uniformly sampled with each new batch, ensuring that the neural network encounters a variety of effective context lengths $\in \bigl[1,\ C_{max}\bigr]$, where $C_{max}$ is the set maximum context length. Because of this, TimesFM learns to handle forecasting based on data with little to no history as well as a complete context.

\subsubsection{Decoder layers with causal attention}

The core of the TimesFM model comprises stacked Transformer decoder layers that operate on the sequence of patches. Each full layer consists of input normalization, followed by a multi-head self-attention and a feed-forward network, both of which are wrapped in a residual connection.
With the model in causal decoder-only mode, the outputs it produces for each patch in the sequence are based solely on the current and preceding tokens, without access to latter parts of the input. This is, yet again, similar to the GPT architecture, and is crucial for a time series forecasting application, as the succeeding tokens in a temporal environment represent events from the future. The causality is enforced in the self-attention layers by masking out all future patches. The output of the decoder can be denoted as

\[ o_i = \mathrm{Decoder}((t_1, \dot{m}_1),\ \ldots,\ (t_i, \dot{m_i}))\]

where:
\begin{itemize}
    \item $ i \in \{1,\ 2,\ \ldots,\ N\}$
    \item $\dot{m_i} = \mathrm{min}\{\tilde{m}_i\}.$ 
\end{itemize}

\subsubsection{Long output patches}

TimesFM produces forecasts in an autoregressive fashion, predicting one patch forward at a time. While it has been shown that for long-term forecasting, predicting all future time steps with a single forecast yields better results than doing so autoregressively with multiple iterations, this is not possible for a foundation model that is meant for variable horizon forecasting. To combat the suboptimal conditions of the task environment, the model is allowed to output patches of data that are longer than their input counterparts. Because of this, the number of autoregressive steps that need to be taken is greatly reduced, thus lessening the negative impact of multiple iterations. The output patch length has one data-driven constraint: it cannot be so long as to exceed the input sequences' entire lengths, which could lead to a collapse in performance. This becomes an issue mainly for time series of high granularity, such as monthly or yearly data. The increased output patch length is achieved by appending another residual block after the decoder layers. It is similar in design to the one responsible for input preprocessing.

\begin{figure}[!h]
		\centering \includegraphics[width=\linewidth]{img/timesfm.jpg}
	\caption{Original TimesFM architecture, reproduced from \cite{timesfm_paper}.}\label{fig:timesfm}
\end{figure}
% \begin{itemize}

%     \item \textbf{Data patching} --- During training, the input times series data is split into fixed-length patches. This is done in order to mimic the tokenization present in large language models. An earlier work [NNSK22] has shown data patching leads to performance improvements in long-term forecasting. Furthermore, data partitioning speeds up inference as the model behaves autoregressively and sees less datapoints per each inference step.

%     \item \textbf{Randomized masking} --- A significant challenge in designing a foundation model is the need to make it work properly at any input length. With constant size inputs fed to the model during training, a notable drop in performance could be observed for time series of any length other than a multiple of the predefined patch length. To combat this, a mask of randomly chosen size is applied to the beginning of each input during training. This can block out parts of the initial patch or even entire patches, effectively feeding inputs of all possible lengths from 1 up to the entire sequence length into the model, thus leading to greater adaptability and uniform performance regardless of sequence length.

% \end{itemize}

\newpage
\subsection{Pretraining and version scope}
The original TimesFM paper describes real and synthetic pretraining data, including web activity and M4 series\cite{timesfm_paper}. The TimesFM 2.5 model card lists GiftEvalPretrain, Wikimedia Pageviews (cutoff November 2023), Google Trends queries (cutoff end of 2022), and synthetic/augmented data\cite{timesfm_card}. This does not provide an auditable list of every training window. No specific mixture percentages, station counts or frequency-dependent context recipe are established here for the executed checkpoint.

The present project loads pretrained weights; it does not recreate pretraining or fine-tune the model. Inference configuration and weight adaptation are distinguished in Section~\ref{sec:sensitivity}, and the limits of benchmark-overlap evidence are discussed in Section~\ref{sec:limitations}.
